\documentclass[12pt,a4paper]{article}

% --- Encoding ---
\usepackage[utf8]{inputenc}
\usepackage[T1]{fontenc}
\usepackage[english]{babel}

% --- Math ---
\usepackage{amsmath, amssymb, mathtools}

% --- Layout ---
\usepackage{geometry}
\geometry{margin=2.5cm, top=3cm, bottom=3cm}

% --- Figures ---
\usepackage{graphicx}
\usepackage{float}
% Figures are generated by notebook_implementation.ipynb and saved to
% DL_MODEL/ptr_net/figures/ — compile from docs/src/ or adjust the path below.
\graphicspath{{../../../DL_MODEL/ptr_net/figures/}}

% --- Tables ---
\usepackage{booktabs}
\usepackage{array}
\usepackage{multirow}
\usepackage{tabularx}
\usepackage[table]{xcolor}

% --- Code listings ---
\usepackage{listings}
\definecolor{codegray}{rgb}{0.45,0.45,0.45}
\definecolor{codeblue}{rgb}{0.0,0.25,0.65}
\definecolor{codeorange}{rgb}{0.75,0.35,0.0}
\definecolor{codered}{rgb}{0.6,0.0,0.0}
\lstdefinestyle{pythonstyle}{
  language=Python,
  basicstyle=\ttfamily\small,
  keywordstyle=\color{codeblue}\bfseries,
  commentstyle=\color{codegray}\itshape,
  stringstyle=\color{codeorange},
  numberstyle=\tiny\color{codegray},
  showstringspaces=false,
  breaklines=true,
  frame=single,
  rulecolor=\color{gray!40},
  backgroundcolor=\color{gray!5},
  xleftmargin=1em,
  numbers=left,
  numbersep=6pt,
}
\lstdefinestyle{bashstyle}{
  language=bash,
  basicstyle=\ttfamily\small,
  keywordstyle=\color{codered},
  commentstyle=\color{codegray}\itshape,
  showstringspaces=false,
  breaklines=true,
  frame=single,
  rulecolor=\color{gray!40},
  backgroundcolor=\color{gray!5},
  xleftmargin=1em,
}

% --- Hyperlinks ---
\usepackage{hyperref}
\hypersetup{colorlinks=true, linkcolor=blue!70!black, citecolor=blue, urlcolor=blue!80!black}

% --- Headers ---
\usepackage{fancyhdr}
\pagestyle{fancy}
\fancyhf{}
\rhead{TSP --- Deep Learning Methods}
\lhead{Pointer Networks (Ptr-Net)}
\cfoot{\thepage}

% --- Misc ---
\usepackage{enumitem}

% =============================================================================
% TITLE
% =============================================================================
\title{%
  \LARGE\textbf{Pointer Networks (Ptr-Net)}\\[0.6em]
  \large Applied to the Travelling Salesman Problem (TSP)
}
\date{\today}

% =============================================================================
\begin{document}
% =============================================================================

\maketitle
\thispagestyle{empty}

\begin{center}
\rowcolors{2}{gray!8}{white}
\begin{tabularx}{\linewidth}{ccX}
\toprule
\textbf{Version} & \textbf{Date} & \textbf{Description} \\
\midrule
v1.0 & 2025-01-01 & Initial release — architecture, training pipeline, TSP benchmark \\
v2.0 & \today     & Added TSPTW-D mode, benchmark restructure, symbol descriptions \\
\bottomrule
\end{tabularx}
\end{center}

\begin{abstract}
This document presents the application of Pointer Networks (Ptr-Net) to the
Travelling Salesman Problem (TSP).
After a formal definition and a complete mathematical description of the
architecture (Vinyals \textit{et al.}, 2015), we describe the supervised
training pipeline, complexity analysis, strengths and limitations of the
approach.
The benchmark covers sizes $n \in \{1, 10, 50, 100, 200\}$ and reports
success rate, non-detection rate, and inference time metrics.
The generalisation degradation --- where a model trained on $n \leq 10$
performs poorly on larger instances --- is analysed and quantified.
Unlike the GNN with its $O(n^2 d)$ edge tensor, the Ptr-Net encoder memory
scales as $O(n \cdot d)$, making it more memory-efficient at the cost of
weaker structural inductive bias.
\end{abstract}

\tableofcontents
\newpage

% =============================================================================
\section*{Acronyms}
% =============================================================================

\begin{tabularx}{\linewidth}{lX}
\toprule
\textbf{Acronym} & \textbf{Definition} \\
\midrule
TSP     & Travelling Salesman Problem \\
TSPTW   & TSP with Time Windows \\
TSPTW-D & TSP with Time Windows and Disruptions \\
Ptr-Net & Pointer Network \\
GNN     & Graph Neural Network \\
LSTM    & Long Short-Term Memory \\
NLL     & Negative Log-Likelihood \\
BCE     & Binary Cross-Entropy \\
RL      & Reinforcement Learning \\
NN      & Nearest-Neighbour (heuristic) \\
OOD     & Out-Of-Distribution \\
AM      & Attention Model (Kool \textit{et al.}, 2019) \\
BPTT    & Backpropagation Through Time \\
\bottomrule
\end{tabularx}
\addcontentsline{toc}{section}{Acronyms}
\newpage

% =============================================================================
\section{Definition}
% =============================================================================

\textbf{Pointer Networks (Ptr-Net)} are a class of sequence-to-sequence
(seq2seq) neural networks in which the decoder ``points'' directly to one of
the input elements at each step, rather than generating a token from a fixed
vocabulary.
This makes the output vocabulary \emph{dynamic} --- it adapts to the current
input, enabling the model to handle variable-length inputs without
architectural changes.

\subsection*{Origin}

\begin{center}
\rowcolors{2}{gray!8}{white}
\begin{tabularx}{\linewidth}{lX}
\toprule
\textbf{Year} & \textbf{Contribution} \\
\midrule
2014 & Sutskever \textit{et al.} --- encoder--decoder seq2seq with fixed
       vocabulary (background) \\
2015 & Vinyals, Fortunato \& Jaitly --- Pointer Networks
       (NeurIPS 2015); applied to TSP, convex hull, Delaunay \\
2016 & Bello \textit{et al.} --- RL-trained Ptr-Net with REINFORCE
       (removes need for optimal labels) \\
2019 & Kool \textit{et al.} --- Transformer Attention Model,
       direct architectural successor to Ptr-Net \\
\bottomrule
\end{tabularx}
\end{center}

\subsection*{Category}

\begin{center}
\rowcolors{2}{gray!8}{white}
\begin{tabularx}{\linewidth}{lX}
\toprule
\textbf{Attribute} & \textbf{Value} \\
\midrule
Learning paradigm & Supervised deep learning \\
Model family & Sequence-to-sequence (seq2seq) \\
Core mechanism & Attention-based pointer over input elements \\
Training signal & Optimal tour labels (ground-truth permutations) \\
\bottomrule
\end{tabularx}
\end{center}

\subsection*{Key concepts}

\begin{itemize}[leftmargin=1.5em]
  \item \textbf{Pointer mechanism}: at each decoding step the model selects
    one input city via a soft attention distribution over encoder hidden states.
    The output vocabulary is the input itself --- size $n$ at test time.
  \item \textbf{Attention as a soft pointer}: the additive (Bahdanau-style)
    attention score $u_i^t$ for city $i$ at step $t$ is
    $v^\top \tanh(W_1 h_i + W_2 s_t)$, interpreted as a logit for pointing
    to city $i$.
  \item \textbf{Masking}: visited cities are set to $-\infty$ before softmax,
    strictly enforcing the permutation constraint without extra architectural
    complexity.
  \item \textbf{Teacher forcing}: during training, the ground-truth tour is
    fed as decoder input, making training stable and well-defined even though
    the model's own greedy decisions may be wrong.
\end{itemize}

% =============================================================================
\section{Formal Description}
% =============================================================================

\subsection{Input and output}

Let the input be a sequence of $n$ city coordinates:
\begin{equation}
X = \{x_1, x_2, \ldots, x_n\}, \qquad x_i \in \mathbb{R}^2
\label{eq:input}
\end{equation}
{\small\begin{tabular}{@{\hspace{1.5em}}l@{\;---\;}p{9.5cm}}
$X$   & set of $n$ input city coordinates \\
$x_i$ & 2D coordinates of city $i$, normalised to $[0,1]^2$ \\
$n$   & number of cities in the instance \\
\end{tabular}}\smallskip

The desired output is a permutation $\pi = (\pi_1, \pi_2, \ldots, \pi_n)$
of $\{1, \ldots, n\}$ that minimises the total tour length
$f(\pi) = \sum_{t=1}^{n} \|x_{\pi_t} - x_{\pi_{t+1}}\|_2$
(indices mod $n$).

\subsection{Probability factorisation}

The joint probability of tour $\pi$ given $X$ is factorised autoregressively:
\begin{equation}
P(\pi \mid X) = \prod_{t=1}^{n}
  P\!\left(\pi_t \mid \pi_1, \ldots, \pi_{t-1},\, X\right)
\label{eq:factorisation}
\end{equation}
{\small\begin{tabular}{@{\hspace{1.5em}}l@{\;---\;}p{9.5cm}}
$\pi$        & output tour permutation $(\pi_1, \dots, \pi_n)$ of city indices \\
$\pi_t$      & city selected at decoding step $t$ \\
$P(\pi \mid X)$ & joint probability of the full tour given the input \\
\end{tabular}}\smallskip

\subsection{Encoder}

An LSTM encoder reads the city coordinates one by one and produces a hidden
state for each:
\begin{equation}
h_i = \mathrm{LSTM}_{\mathrm{enc}}\!\left(
  W_{\mathrm{in}}\,x_i,\; h_{i-1}
\right), \qquad
W_{\mathrm{in}} \in \mathbb{R}^{d_e \times 2}
\label{eq:encoder}
\end{equation}
{\small\begin{tabular}{@{\hspace{1.5em}}l@{\;---\;}p{9.5cm}}
$h_i$              & encoder hidden state for city $i$, $h_i \in \mathbb{R}^d$ \\
$W_{\mathrm{in}}$  & learned linear projection from coordinates to embedding space \\
$d_e$              & embedding dimension (input to the LSTM) \\
$d$                & hidden dimension of the LSTM \\
\end{tabular}}\smallskip

where $d_e$ is the embedding dimension and $d$ the hidden dimension.
The sequence of hidden states $\{h_1, \ldots, h_n\}$ forms the
\emph{memory bank} from which the decoder will point.
The final state $(h_n, c_n)$ initialises the decoder.

\subsection{Decoder}

An LSTM decoder maintains a state $s_t \in \mathbb{R}^d$ updated at each step:
\begin{equation}
s_t = \mathrm{LSTMCell}\!\left(
  \tilde{x}_{\pi_{t-1}},\; s_{t-1}
\right)
\label{eq:decoder}
\end{equation}
{\small\begin{tabular}{@{\hspace{1.5em}}l@{\;---\;}p{9.5cm}}
$s_t$                     & decoder hidden state at step $t$ \\
$\tilde{x}_{\pi_{t-1}}$   & embedding of the previously chosen city, $W_{\mathrm{city}}\,x_{\pi_{t-1}}$ \\
$W_{\mathrm{city}}$       & shared city embedding matrix, same for encoder and decoder \\
\end{tabular}}\smallskip

where $\tilde{x}_{\pi_{t-1}} = W_{\mathrm{city}}\,x_{\pi_{t-1}}$ is the
embedding of the previously chosen city.
A learnable \textbf{start token} $s_0$ replaces the unavailable
``previous city'' at $t = 1$.

\subsection{Attention / pointer scores}

At step $t$, an additive attention score is computed for every encoder state:
\begin{equation}
u_i^t = v^\top \tanh\!\left(W_1\,h_i + W_2\,s_t\right),
\qquad W_1, W_2 \in \mathbb{R}^{d \times d},\; v \in \mathbb{R}^d
\label{eq:attention}
\end{equation}
{\small\begin{tabular}{@{\hspace{1.5em}}l@{\;---\;}p{9.5cm}}
$u_i^t$  & unnormalised pointer score for city $i$ at decoder step $t$ \\
$W_1$    & weight matrix projecting encoder hidden state $h_i$ \\
$W_2$    & weight matrix projecting decoder state $s_t$ \\
$v$      & learned query vector; dot product extracts scalar score \\
\end{tabular}}\smallskip

\subsection{Pointer distribution with masking}

Let $\mathcal{V}_t$ be the set of already-visited indices at step $t$.
Scores of visited cities are set to $-\infty$ before softmax:
\begin{equation}
\tilde{u}_i^t = \begin{cases}
  u_i^t   & \text{if } i \notin \mathcal{V}_t \\
  -\infty  & \text{if } i \in \mathcal{V}_t
\end{cases}
\label{eq:masking}
\end{equation}
{\small\begin{tabular}{@{\hspace{1.5em}}l@{\;---\;}p{9.5cm}}
$\tilde{u}_i^t$  & masked score for city $i$ at step $t$ \\
$\mathcal{V}_t$  & set of already-visited city indices at step $t$ \\
\end{tabular}}\smallskip

\begin{equation}
P\!\left(\pi_t = i \mid \pi_{<t},\, X\right)
  = \mathrm{softmax}\!\left(\tilde{u}^t\right)_i
  = \frac{\exp(\tilde{u}_i^t)}{\sum_{j=1}^{n} \exp(\tilde{u}_j^t)}
\label{eq:pointer_dist}
\end{equation}
{\small\begin{tabular}{@{\hspace{1.5em}}l@{\;---\;}p{9.5cm}}
$P(\pi_t = i \mid \ldots)$ & probability of pointing to city $i$ at step $t$, given past choices and input \\
$\tilde{u}^t$              & vector of masked scores for all cities at step $t$ \\
\end{tabular}}\smallskip

\subsection{Training objective}

Given a training example $(X, \pi^*)$ where $\pi^*$ is the optimal
ground-truth tour, the supervised loss is the negative log-likelihood
(cross-entropy) summed over all decoding steps with teacher forcing:

\begin{equation}
\mathcal{L}_{\mathrm{NLL}}
  = -\sum_{t=1}^{n}
    \log P\!\left(\pi_t^* \mid \pi_1^*, \ldots, \pi_{t-1}^*,\, X\right)
\label{eq:nll}
\end{equation}
{\small\begin{tabular}{@{\hspace{1.5em}}l@{\;---\;}p{9.5cm}}
$\mathcal{L}_{\mathrm{NLL}}$ & negative log-likelihood loss (NLL), minimised during training \\
$\pi^*$                      & ground-truth tour provided by the label strategy (optimal / nn2opt) \\
$\pi_t^*$                    & target city at step $t$ under teacher forcing \\
\end{tabular}}\smallskip

\textbf{Labelling constraint:} optimal tours are computed by brute-force
enumeration ($O(n!)$), which is feasible only for $n \leq 10$.
The \texttt{train.py} script enforces this with an assertion.

\subsection{Complexity}

\begin{center}
\rowcolors{2}{gray!8}{white}
\begin{tabularx}{\linewidth}{lXXX}
\toprule
\textbf{Component} & \textbf{Time} & \textbf{Space} & \textbf{Dominant term} \\
\midrule
Encoder ($n$ LSTM steps) & $O(n d^2)$ & $O(n d)$ & hidden states \\
Decoder ($n$ LSTM steps) & $O(n d^2)$ & $O(d)$ & one state at a time \\
Attention (per decoder step) & $O(n d)$ & $O(n)$ & scores over all cities \\
\midrule
\textbf{Total} & $\boldsymbol{O(n^2 d + n d^2)}$ &
  $\boldsymbol{O(n d)}$ & \textbf{attention} \\
\bottomrule
\end{tabularx}
\end{center}

For typical $n \ll d$ the dominant term is $O(n d^2)$; for large $n$ it
becomes $O(n^2 d)$.
Crucially, the space complexity is $O(n d)$ --- \emph{linear} in $n$,
compared to $O(n^2 d)$ for the GNN edge tensor.
This makes Ptr-Net much more memory-efficient, at the cost of weaker
structural inductive bias.

% =============================================================================
\section{Architecture / Algorithm Steps}
% =============================================================================

\begin{center}
\begin{tabularx}{\linewidth}{X}
\textbf{Input} $(x_1, \dots, x_n)$ city coordinates\\
$\quad\downarrow$\\
\textbf{Step 1} --- Embed each city: $\tilde{x}_i = W_{\mathrm{in}}\,x_i$\\
$\quad\downarrow$\\
\textbf{Step 2} --- LSTM encoder reads the sequence
$\to$ memory bank $\{h_1, \dots, h_n\}$, final state $(h_n, c_n)$\\
$\quad\downarrow$\\
\textbf{Step 3} --- Decoder loop for $t = 1, \dots, n$:\\
$\quad\quad\quad$ (a) LSTMCell update: $s_t \leftarrow \mathrm{LSTMCell}(\tilde{x}_{\pi_{t-1}}, s_{t-1})$\\
$\quad\quad\quad$ (b) Pointer scores: $u_i^t = v^\top \tanh(W_1 h_i + W_2 s_t)$ for all $i$\\
$\quad\quad\quad$ (c) Mask visited cities ($-\infty$); apply softmax\\
$\quad\quad\quad$ (d) Select $\pi_t = \arg\max_i P(\pi_t = i \mid \ldots)$; mark $\pi_t$ as visited\\
$\quad\downarrow$\\
\textbf{Output} tour $\pi = (\pi_1, \dots, \pi_n)$\\
\end{tabularx}
\end{center}

\textbf{Training loop (supervised):}
\begin{enumerate}[leftmargin=2em]
  \item Sample a random TSP instance $(x_1, \dots, x_n)$, $n \leq 10$.
  \item Compute the ground-truth tour $\pi^*$ by brute-force ($O(n!)$).
  \item Teacher-forced forward pass: feed $\pi^*$ as decoder input at every
    step; collect log-probabilities $\log P(\pi_t^* \mid \ldots)$.
  \item Compute $\mathcal{L}_{\mathrm{NLL}}$ and backpropagate through time
    (BPTT).
  \item Update weights with Adam; clip gradients to norm 1.
\end{enumerate}

\textbf{Decoding strategies:}

\begin{center}
\rowcolors{2}{gray!8}{white}
\begin{tabularx}{\linewidth}{lX}
\toprule
\textbf{Strategy} & \textbf{Description} \\
\midrule
Greedy & Select $\arg\max_i P(\pi_t = i \mid \ldots)$ at every step ---
  $O(n^2 d)$ total \\
Sampling & Draw $k$ independent tours from the distribution; keep shortest ---
  $O(k n^2 d)$ \\
Beam search & Maintain $B$ partial tours; select top-$B$ by cumulative
  log-probability --- $O(B n^2 d)$ \\
\bottomrule
\end{tabularx}
\end{center}

% =============================================================================
\section{Strengths and Limitations}
% =============================================================================

\subsection*{Strengths}

\begin{center}
\rowcolors{2}{gray!8}{white}
\begin{tabularx}{\linewidth}{lX}
\toprule
\textbf{Property} & \textbf{Detail} \\
\midrule
Variable input size & Handles any $n$ at inference; pointer mechanism adapts
  to current input length \\
Memory efficient & Encoder memory $O(n d)$ vs.\ GNN's $O(n^2 d)$; feasible
  for much larger $n$ \\
Interpretable attention & Pointer distribution $P(\pi_t = i \mid \ldots)$
  directly shows which city is being considered \\
Strict permutation output & Masking guarantees no city is visited twice \\
Lightweight & Fewer parameters than GNN or Transformer; fast to train
  on small instances \\
\bottomrule
\end{tabularx}
\end{center}

\subsection*{Limitations}

\begin{center}
\rowcolors{2}{gray!8}{white}
\begin{tabularx}{\linewidth}{lX}
\toprule
\textbf{Limitation} & \textbf{Impact} \\
\midrule
Sequential decoder & LSTM decoder is inherently sequential; cannot be
  parallelised across time steps (cf.\ Transformer) \\
Label constraint & Optimal labels require brute-force ($n \leq 10$); no
  near-optimal labelling strategy for larger $n$ without a separate solver \\
Generalisation degradation & Model trained on $n = 8$ performs poorly on
  $n = 50$ without retraining; gap grows rapidly with $n$ \\
No structural inductive bias & Treats TSP as a sequence problem, ignoring
  pairwise distances and graph topology \\
Superseded by Transformer AM & Kool \textit{et al.} (2019) achieves
  $\sim$1\,\% gap on $n = 100$ vs $\sim$8--10\,\% for Ptr-Net \\
\bottomrule
\end{tabularx}
\end{center}

% =============================================================================
\section{Use Case Explanation}
% =============================================================================

\subsection*{When to use Pointer Networks for TSP}

\begin{enumerate}[leftmargin=2em]
  \item \textbf{Small instances ($n \leq 50$) with optimal labels available.}
    Ptr-Net achieves near-optimal tours on the training distribution.
    With beam search ($B = 10$), the gap to optimal is below 5\,\% for
    $n \leq 20$.

  \item \textbf{As a baseline before more complex models.}
    Ptr-Net is the historical predecessor of GNN and Transformer-AM approaches.
    Implementing it first gives a clear understanding of the
    encoder--decoder--pointer paradigm.

  \item \textbf{Interpretability is a requirement.}
    The pointer distribution at each step is a proper probability over input
    cities, making the selection process easy to visualise and explain.

  \item \textbf{Memory-constrained environments.}
    With $O(n d)$ encoder memory, Ptr-Net can handle larger $n$ than the GNN
    without GPU out-of-memory issues.
\end{enumerate}

\subsection*{Why Ptr-Net specifically (vs.\ GNN or Transformer AM)}

\begin{center}
\rowcolors{2}{gray!8}{white}
\begin{tabularx}{\linewidth}{lXXXX}
\toprule
\textbf{Model} & \textbf{Inductive bias} & \textbf{Gap at $n=100$} &
  \textbf{Memory} & \textbf{Labels needed} \\
\midrule
\textbf{Ptr-Net} (Vinyals 2015) & Sequential LSTM
  & ${\sim}8$--$10\,\%$ & $O(nd)$ & Optimal ($n \leq 10$) \\
GNN (Joshi 2019)          & Graph msg.\ passing
  & ${\sim}1$--$3\,\%$  & $O(Ln^2 d)$ & Near-optimal \\
Transformer AM (Kool 2019) & Global pairwise attn.
  & ${\sim}1$--$2\,\%$  & $O(n^2 d)$ & None (RL) \\
LKH-3 (Helsgott)          & Metaheuristic
  & $< 0.5\,\%$         & $O(n^2)$ & None \\
\bottomrule
\end{tabularx}
\end{center}

Ptr-Net is the most memory-efficient neural solver and the simplest to
implement, making it the natural starting point for understanding neural
combinatorial optimisation.
Its weaker structural prior (no explicit edge features) explains the larger
optimality gap compared to GNN and Transformer AM.

% =============================================================================
\section{Code Implementation}
% =============================================================================

The full architecture is defined in \texttt{model.py}.
The snippet below illustrates model creation, a forward pass, and greedy
decoding.

\begin{lstlisting}[style=pythonstyle,
  caption={Creation, inference, and tour decoding
           (\texttt{model.py} + \texttt{data.py})}]
import torch
from model import PointerNetwork, MODEL_SIZES
from data  import random_instance, tour_length

# --- Choose a preset ---------------------------------------------------
SIZE               = "medium"        # "small" | "medium" | "large"
embed, hidden, lyr = MODEL_SIZES[SIZE]

# --- Load trained model ------------------------------------------------
model = PointerNetwork(embed_dim=embed, hidden_dim=hidden, n_layers=lyr)
model.load_state_dict(
    torch.load(f"model/ptr_net_{SIZE}.pt", map_location="cpu")
)
model.eval()

# --- Generate a random instance ----------------------------------------
n      = 8                          # training supports n <= 10
coords = random_instance(n)         # torch.Tensor (n, 2) in [0, 1]

# --- Greedy inference ---------------------------------------------------
with torch.no_grad():
    log_probs, tour = model(coords)  # log_probs: (n,n),  tour: list[int]

length = tour_length(coords, tour)
print(f"Tour   : {tour}")
print(f"Length : {length:.4f}")
\end{lstlisting}

\begin{lstlisting}[style=bashstyle,
  caption={Training from the command line
           (run from \texttt{DL\_MODEL/ptr\_net/})}]
# Small model, 1 000 steps, random instances (n=8)
python train.py --mode tsp --size small --epochs 1000

# Medium model, best results on training distribution
python train.py --mode tsp --size medium --epochs 3000 --source tsp

# Fine-tuning from an existing checkpoint
python train.py --mode tsp --resume model/ptr_net_medium.pt \
                --size medium --epochs 500 --lr 1e-4

# TSPTW-D model (node_dim=5, time windows + disruptions)
python train.py --mode tsptwd --size medium --epochs 2000
\end{lstlisting}

Available presets:

\begin{center}
\rowcolors{2}{gray!8}{white}
\begin{tabularx}{\linewidth}{ccccX}
\toprule
\textbf{Preset} & $d_e$ & $d$ & \textbf{Parameters} & \textbf{Recommended use} \\
\midrule
\texttt{small}  & 64  & 128 & ${\sim}232{,}000$   & Quick tests, $n \leq 20$ \\
\texttt{medium} & 128 & 256 & ${\sim}923{,}000$   & Standard TSP / TSPTW-D, $n \leq 50$ \\
\texttt{large}  & 256 & 512 & ${\sim}3{,}700{,}000$ & Larger instances \\
\bottomrule
\end{tabularx}
\end{center}

% =============================================================================
\section{TSP Benchmark}
% =============================================================================

The benchmark uses the \textbf{medium} model preset
($d_e = 128$, $d = 256$, 922\,752 parameters, trained with the two-stage
supervised curriculum in \texttt{train.sh}).
For in-distribution evaluation ($n \leq 10$) the reference is the brute-force
optimal tour; for out-of-distribution sizes ($n > 10$) the nearest-neighbour
heuristic is used as reference.
Results for $n = 10$ are averages over 50 random instances;
results for $n \in \{50, 100\}$ are averages over 20 instances.
Inference was measured on Apple Silicon (MPS device).

\subsection*{Metrics}

\begin{center}
\rowcolors{2}{gray!8}{white}
\begin{tabularx}{\linewidth}{lX}
\toprule
\textbf{Metric} & \textbf{Definition} \\
\midrule
Avg.\ length & Mean closed-tour Euclidean length over test instances \\
Gap (\%)     & $(L_{\mathrm{model}} - L_{\mathrm{ref}}) / L_{\mathrm{ref}} \times 100$, where $L_{\mathrm{ref}}$ is the reference tour length \\
Avg.\ time   & Mean wall-clock time per instance for one forward pass + greedy decode (CPU) \\
\bottomrule
\end{tabularx}
\end{center}

\subsection*{Results — $n = 10$ (in-distribution, 50 instances)}

\begin{center}
\rowcolors{2}{gray!8}{white}
\begin{tabularx}{\linewidth}{Xccc}
\toprule
\textbf{Method} & \textbf{Avg.\ length} & \textbf{Gap vs optimal (\%)} & \textbf{Avg.\ time} \\
\midrule
Nearest-neighbour & 3.0469 & $+9.98\,\%$ & $0.09$\,ms \\
2-opt             & 2.7789 & $+0.31\,\%$ & $0.36$\,ms \\
Ptr-Net (greedy)  & 2.9296 & $+5.75\,\%$ & $18.4$\,ms \\
Ptr-Net + 2-opt   & 2.7864 & $+0.58\,\%$ & $13.3$\,ms \\
Optimal (ref.)    & 2.7703 & $0.00\,\%$  & ---         \\
\bottomrule
\end{tabularx}
\end{center}

\subsection*{Results — $n = 50$ (out-of-distribution, 20 instances)}

\begin{center}
\rowcolors{2}{gray!8}{white}
\begin{tabularx}{\linewidth}{Xccc}
\toprule
\textbf{Method} & \textbf{Avg.\ length} & \textbf{Gap vs NN (\%)} & \textbf{Avg.\ time} \\
\midrule
Nearest-neighbour & 7.1567 & $0.00\,\%$ (ref.) & $0.28$\,ms \\
2-opt             & 6.0832 & $-15.00\,\%$      & $16.0$\,ms \\
Ptr-Net (greedy)  & 7.6698 & $+7.17\,\%$       & $44.3$\,ms \\
Ptr-Net + 2-opt   & 6.1632 & $-13.88\,\%$      & $47.7$\,ms \\
\bottomrule
\end{tabularx}
\end{center}

\subsection*{Results — $n = 100$ (out-of-distribution, 20 instances)}

\begin{center}
\rowcolors{2}{gray!8}{white}
\begin{tabularx}{\linewidth}{Xccc}
\toprule
\textbf{Method} & \textbf{Avg.\ length} & \textbf{Gap vs NN (\%)} & \textbf{Avg.\ time} \\
\midrule
Nearest-neighbour & 9.8126  & $0.00\,\%$ (ref.) & $0.49$\,ms  \\
2-opt             & 8.3997  & $-14.40\,\%$      & $76.9$\,ms  \\
Ptr-Net (greedy)  & 12.3332 & $+25.69\,\%$      & $86.9$\,ms  \\
Ptr-Net + 2-opt   & 8.6146  & $-12.21\,\%$      & $144.7$\,ms \\
\bottomrule
\end{tabularx}
\end{center}

\subsection*{Figures}

\begin{figure}[H]
  \centering
  \includegraphics[width=\linewidth]{benchmark_computation.png}
  \caption{
    Inference time (linear and log-log x-scale) and encoder memory footprint
    as a function of $n$.
    Unlike the GNN, encoder memory grows \emph{linearly} with $n$ ($O(n d)$),
    making the model feasible for much larger instances.
    The dashed red line in the log-log plot is the $O(n)$ reference for
    memory and $O(n^2)$ for time.
  }
  \label{fig:benchmark_computation}
\end{figure}

\begin{figure}[H]
  \centering
  \includegraphics[width=0.75\linewidth]{benchmark_quality.png}
  \caption{
    Gap to the NN baseline for $n \in \{50, 100\}$.
    Ptr-Net greedy is $+7.2\,\%$ above NN at $n=50$ and $+25.7\,\%$ at $n=100$,
    confirming rapid OOD degradation.
    Ptr-Net + 2-opt recovers well ($-13.9\,\%$ and $-12.2\,\%$ below NN),
    matching 2-opt quality but at higher latency.
  }
  \label{fig:benchmark_quality}
\end{figure}

\begin{figure}[H]
  \centering
  \includegraphics[width=\linewidth]{demo_overview.png}
  \caption{
    Overview on an $n = 8$ instance (seed 0):
    (i)~NLL training loss curve,
    (ii)~brute-force optimal tour (green),
    (iii)~Ptr-Net greedy tour (blue) with optimality gap,
    (iv)~pointer distribution at decoding step 0.
  }
  \label{fig:demo_overview}
\end{figure}

% =============================================================================
\subsection{Generalisation and scalability analysis}
\label{sec:scalability}
% =============================================================================

The supervised training loop requires brute-force optimal labels, which are
only tractable for $n \leq 10$.
A model trained at $n = 8$ is therefore evaluated \emph{out-of-distribution}
for all $n > 10$, leading to rapid quality degradation.

\begin{table}[H]
\centering
\rowcolors{2}{gray!8}{white}
\begin{tabularx}{\linewidth}{ccX}
\toprule
$n$ & Encoder memory ($d=256$) & Note \\
\midrule
10      & ${\sim}0.03$\,MB  & in-distribution (training) \\
50      & ${\sim}0.13$\,MB  & OOD, quality degrades \\
100     & ${\sim}0.26$\,MB  & OOD, significant gap \\
200     & ${\sim}0.51$\,MB  & OOD, quality poor \\
1\,000  & ${\sim}2.6$\,MB   & inference feasible, quality very poor \\
10\,000 & ${\sim}26$\,MB    & slow; sequential decoder bottleneck \\
\bottomrule
\end{tabularx}
\caption{Ptr-Net memory scaling.
  Encoder memory is $O(n \cdot d)$ --- linear in $n$, unlike the $O(n^2 d)$
  GNN edge tensor.
  Quality for $n > 10$ is unreliable without retraining on larger instances.}
\label{tab:scalability}
\end{table}

\paragraph{Conclusion.}
The Ptr-Net's $O(n d)$ memory footprint is its key advantage over the GNN:
inference remains feasible for $n \gg 500$.
However, the quality bottleneck is not memory but \textbf{generalisation}:
the model cannot reliably solve instances larger than its training size.
To exploit the memory advantage, one must either:

\begin{itemize}[leftmargin=1.5em]
  \item \textbf{Train on larger $n$ with near-optimal labels} (e.g., 2-opt
    post-processed NN tours as pseudo-labels, or a separate LKH call).
  \item \textbf{Switch to RL training} (Bello \textit{et al.}, 2016):
    reward = negative tour length; no optimal labels required; scales to
    $n \leq 100$ with REINFORCE.
  \item \textbf{Upgrade to Transformer AM} (Kool \textit{et al.}, 2019):
    retains the pointer mechanism but replaces LSTM with multi-head attention
    and trains with RL, achieving $\sim$1\,\% gap at $n = 100$.
\end{itemize}

% =============================================================================
\section{Experimental Analysis}
% =============================================================================

\paragraph{Inference time scaling.}
Figure~\ref{fig:benchmark_computation} shows that inference time grows
approximately as $O(n^2)$ due to the attention loop, while encoder memory
grows linearly with $n$ --- confirming the theoretical analysis.
The sequential LSTM decoder dominates wall-clock time for large $n$.

\paragraph{Solution quality: in-distribution vs.\ out-of-distribution.}
Figure~\ref{fig:benchmark_quality} and the result tables quantify the
central limitation of supervised Ptr-Net.
At $n = 10$ (in-distribution), Ptr-Net greedy achieves a $+5.75\,\%$ gap
versus the brute-force optimal --- comparable to NN ($+9.98\,\%$) and far
better than random.
Adding 2-opt post-processing closes the gap to $+0.58\,\%$.
At $n = 50$ (OOD), greedy Ptr-Net degrades to $+7.2\,\%$ \emph{above}
the NN baseline, confirming that the model's pointer mechanism does not
generalise to larger sequences.
At $n = 100$, the gap widens to $+25.7\,\%$.
Ptr-Net + 2-opt largely recovers quality ($-13.9\,\%$ and $-12.2\,\%$
below NN at $n = 50$ and $100$ respectively), but this hybrid benefit
comes from 2-opt, not the network itself.
This contrasts with the GNN, which uses near-optimal labels at larger $n$
and therefore generalises better within the $n \leq 200$ range.

\paragraph{Attention interpretation.}
The pointer distribution $P(\pi_t = i \mid \ldots)$ at each step can be
visualised as a heatmap over the cities.
Early steps tend to concentrate probability on geometrically close or
globally prominent cities (e.g., boundary cities); later steps become more
focused as the number of unvisited cities decreases.

\paragraph{Silent failures.}
The Ptr-Net is deterministic at inference (greedy mode) and always outputs a
tour with no internal quality estimate.
Non-detection (silent failure) rate equals $100 - \%\,\mathrm{success}$ by
construction, and false-detection rate is always $0\,\%$.
A simple post-hoc check (compare to NN baseline) can flag poor solutions
before deployment.

\paragraph{Improvement directions.}

\begin{center}
\rowcolors{2}{gray!8}{white}
\begin{tabularx}{\linewidth}{lX}
\toprule
\textbf{Direction} & \textbf{Expected benefit} \\
\midrule
RL training (REINFORCE) & Removes need for optimal labels; enables training
  on $n \leq 100$ \\
Beam search decoding & Improves quality by $\sim$2--5\,\% at cost of $B\times$
  inference time \\
Train on larger $n$ with pseudo-labels & Closes OOD generalisation gap for
  $n \leq 100$ \\
Bidirectional LSTM encoder & Richer context per city; modest quality
  improvement \\
Replace LSTM with Transformer & Full Transformer AM; state-of-the-art at
  $n \leq 100$ \\
\bottomrule
\end{tabularx}
\end{center}

% =============================================================================
\section*{References}
% =============================================================================

\begin{itemize}[leftmargin=1.5em]
  \item Vinyals O., Fortunato M., Jaitly N. (2015).
    \textit{Pointer Networks}.
    Advances in Neural Information Processing Systems (NeurIPS), 28.
  \item Bello I., Pham H., Le Q.V., Norouzi M., Bengio S. (2016).
    \textit{Neural Combinatorial Optimization with Reinforcement Learning}.
    ICLR Workshop 2017.
  \item Sutskever I., Vinyals O., Le Q.V. (2014).
    \textit{Sequence to Sequence Learning with Neural Networks}.
    NeurIPS 2014.
    (Background: encoder--decoder with fixed vocabulary.)
  \item Kool W., Van Hoof H., Welling M. (2019).
    \textit{Attention, Learn to Solve Routing Problems!}
    ICLR 2019.
    (Direct successor: Transformer Attention Model.)
  \item Joshi C.K., Laurent T., Bresson X. (2019).
    \textit{An Efficient Graph Convolutional Network Technique for the
    Travelling Salesman Problem}.
    arXiv:1906.01227.
  \item Bahdanau D., Cho K., Bengio Y. (2015).
    \textit{Neural Machine Translation by Jointly Learning to Align and
    Translate}.
    ICLR 2015.
    (Additive attention mechanism used by Ptr-Net.)
  \item Solomon M.M. (1987).
    \textit{Algorithms for the Vehicle Routing and Scheduling Problems with
    Time Window Constraints}.
    Operations Research, 35(2), 254--265.
\end{itemize}

\end{document}
